\documentclass[11pt,a4paper]{article}
\usepackage[utf8]{inputenc}
\usepackage[left=2cm,right=2cm,text={18cm,24cm},top=2cm]{geometry}
\usepackage[english]{babel}
\usepackage{graphicx}
\usepackage{verbatim}
\renewcommand{\familydefault}{\sfdefault}


\begin{document}
    \begin{titlepage}
        \begin{center}
            % FIT logo
            \includegraphics[scale=0.65]{include/fit.pdf} \\

            \LARGE{
                \textbf{
                    Project documentation} \\
                Compiler for the IFJ22 language} \\
                
            \vspace{2cm}
            
            \Large{
                Team xstrel03 \\
                Variant TRP \\
            }
                
            \vspace{2cm}
            
            \normalsize{}
            \today{}

            \vspace{2cm}

            \begin{tabular}{l l l}
                \textbf{Matyáš Strelec} & \textbf{(xstrel03)}   & \quad X\% \\
                Ondřej Seidl            & (xseidl06)            & \quad X\% \\
                Maxmilián Nový          & (xnovym00)            & \quad X\% \\
                Dominik Klon            & (xklond00)            & \quad X\% \\
            \end{tabular}
        \end{center}
    \end{titlepage}

    \tableofcontents

    \newpage{}

    \section{Introduction}
    V dokumentaci popisujte návrh (části překladače a předávání informací mezi nimi),
    implementaci (použité datové struktury, tabulku symbolů, generování kódu),
    vývojový cyk- lus, způsob práce v týmu, speciální použité techniky a algoritmy
    a různé odchylky od před- nášené látky či tradičních přístupů. Nezapomínejte také
    citovat literaturu a uvádět refe- rence na čerpané zdroje včetně správné citace
    převzatých částí (obrázky, magické konstanty, vzorce). Nepopisujte záležitosti
    obecně známé či přednášené na naší fakultě.

    \section{Division of work}
    Rozdělení práce mezi členy týmu (uveďte kdo a jak se podílel na jednotlivých
    částech projektu; povinně zdůvodněte odchylky od rovnoměrného rozdělení bodů).

    \section{Lexical analysis}
    Diagram konečného automatu, který specifikuje lexikální analyzátor.

    \section{Syntactic analysis}
    LL-gramatiku,LL-tabulku a precedenční tabulku, podle kterých jste implementovali
    váš syntaktický analyzátor.

    \section{Implementation details}
    Stručný popis členění implementačního řešení včetně názvů souborů, kde jsou jednotlivé
    části včetně povinných implementovaných metod překladače k nalezení.

\end{document}
